
\documentclass[12pt]{article}
\usepackage[margin=1in]{geometry}
\usepackage{hyperref}
\begin{document}

\section*{Introduction: Motivation and Epistemological Commitments}

The Scalar--Angular--Torsion (SAT) framework began not as a fully formed theory, but as a simple act of taking existing representations literally. In particular, the idea emerged from treating the familiar practice of plotting particle worldlines through spacetime as more than a mathematical convenience---as a physical ontology. If a particle's path through spacetime is not just a trace of past events but an actual physical structure, then the geometry of those worldlines---and their interaction with an evolving time surface---may encode the dynamical content of the physical world.

From this initial move, a sequence of geometric and structural relationships began to emerge. The intersection of worldlines with a propagating time surface suggested a new lens on mass and energy; the helical forms of filament paths hinted at internal symmetries; bundles and knots implied a topological basis for particle classes. As these ideas developed, they began to reflect---and in some cases reproduce---the structure of existing physical theories. General Relativity, the Standard Model, and elements of quantum field theory all began to emerge from the behavior and statistics of these minimal geometric objects: filaments and their interactions with the time surface.

The SAT program is not an attempt to displace or supersede these well-established frameworks. Quite the opposite: the Standard Model, General Relativity, and Quantum Mechanics are taken here as foundational constraints. They are the ``tent poles'' of physical theory---successful, empirically grounded, and indispensable. SAT's goal is to offer a minimal geometric and topological substrate from which these theories might emerge, not through imposed axioms or field definitions, but as natural consequences of deeper structural relationships.

This approach commits itself to a core methodological principle: \emph{the data is primary}. Observed phenomena---particle properties, interaction strengths, cosmological parameters---must be taken as fixed inputs. The role of theory is not to overwrite them, but to organize, relate, and explain them. Accordingly, SAT was built not to predict the known, but to account for it in a way that is compact, cohesive, and geometrically grounded.

And yet, when a model constructed to mirror existing relationships begins to generate \emph{unexpected internal coherence}, or reproduces results it was not explicitly designed to fit, it crosses a threshold. These are not mere coincidences or curve-fits, but structural reflections---points where geometry and empirical fact align without being manually forced to do so. SAT, we believe, sits at such a threshold. The theory was not tuned to reproduce gauge coupling values, yet it does. It was not constructed with hadron masses in mind, yet they follow from its topological mass scaling. It was not invented to model inflation, yet its geometry naturally produces early-time acceleration and a viable spectrum.

This is the space in which SAT operates---not yet a complete or final theory, but a framework that maps known physics to intuitive geometric principles, and in doing so, offers a unifying structure with surprising reach. The work that follows is an attempt to formalize those intuitions, clarify their mathematical structure, and lay out both the successes and the limitations of the current model.

\end{document}
